\subsection{正规空间的定义}

关于可一致化空间的公理 \((\mathbf{O}_{\mathbf{IV}})\) （§ 1，第 5 号）可叙述成如下形式：给定任一闭集 A 与任一点 \(x \in \complement \mathbf{A}\) ，存在 \(\mathbf{X}\) 到 {[}0, 1{]} 中的连续映射，它在 x 处等于 o 而在 A 的每一点处等于 i；这个性质又可表述为：在可一致化空间中，可用连续实值函数分离一点与一个闭集（不含该点的）。

我们现在来研究这样的空间：在其中同样可以用连续实值函数分离两个不相交的闭集：

定义 1. 拓扑空间 X 称为正规的，如果它是豪斯多夫的且满足下列公理：

\((\mathbf{O}_{\mathbf{v}})\) 若 A 与 B 是 X 的任意两个不相交闭子集，则存在 X 到 {[}o, i{]} 中的连续映射，它在 A 的每一点处等于 o 而在 B 的每一点处等于 i。

显然每个正规空间都是完全正则的；但存在不是正规的完全正则空间（见习题 9、10、13、26 及 § 5 习题 15 与 16）。

公理 \((\mathrm{O}_{\mathrm{V}})\) 的叙述像 \((\mathrm{O}_{\mathrm{IV}})\) 的一样，涉及实直线 R 作为辅助集。但有一个与 \((\mathrm{O}_{\mathrm{V}})\) 等价且不涉及任何辅助集的条件：

定理 1（乌雷松）。公理 \((\mathrm{O}_{\mathrm{V}})\) 等价于下列公理：

\((\mathbf{O}_{\mathbf{V}}^{\prime})\) 若 \(\mathbf{A}\) 与 \(\mathbf{B}\) 是 \(\mathbf{X}\) 的任意两个不相交闭子集，则存在两个不相交的开集 \(\mathbf{U}, \mathbf{V}\) 使得 \(\mathbf{A} \subset \mathbf{U}\) 且 \(\mathbf{B} \subset \mathbf{V}\) 。

立即得出 \((\mathrm{O}_{\mathrm{V}})\) 蕴含 \((\mathrm{O}_{\mathrm{V}}^{\prime})\) ，因为若 f 是 X 到 {[}0, I{]} 中的连续映射，它在 A 上等于 o 而在 B 上等于 I，则开集 \(\overline{f}^{1}([0, I/2[)\text{ 与 }\overline{f}^{1}([I/2,I])\text{ 分别包含 A 与 B，且互不相交。}\)

为证明其逆，首先注意 \((\mathrm{O}_{\mathrm{V}}^{\prime})\) 等价于下列公理：

\((\mathbf{O}_{\mathbf{V}}^{\prime \prime})\) 给定任一闭集 \(\mathbf{A}\) 与任一开邻域 \(\mathbf{V}\) （\(\mathbf{A}\) 的），存在开邻域 \(\mathbf{W}\) （\(\mathbf{A}\) 的）使得 \(\overline{\mathbf{W}} \subset \mathbf{V}\) 。

若存在连续映射 \(f: \mathbf{X} \to [-\mathbf{i}, +\mathbf{i}]\) ，它等于 \(-\mathbf{i}\) 于 \(\mathbf{A}\) 上而等于 \(+\mathbf{i}\) 于 \(\mathbf{B}\) 上，并且令 \(\mathbf{U}(t) = \overline{f}^1([-\mathbf{i}, t[)\) （对每个 \(t \in [\mathbf{o}, \mathbf{i}]\) ），则我们就定义了 \(\mathbf{X}\) 中的一族开集，以 \([\mathbf{o}, \mathbf{i}]\) 为指标集，使得 (i) \(\mathbf{A} \subset \mathbf{U}(\mathbf{o})\) ，(ii) \(\mathbf{B} \subset \mathbf{C}\mathbf{U}(\mathbf{i})\) ，且 (iii) 对每对实数 \(t, t'\) （满足 \(\mathbf{o} \leqslant t < t' \leqslant \mathbf{i}\) ）有

\[
\overline {{{\mathbf {U}}}} (t) \subset \mathbf {U} (t ^ {\prime});\tag{1}
\]

因为 \(\mathbf{U}(t)\) 含于闭集 \(\overline{f}([-\mathrm{i}, t])\) 。反之，设我们已定义了具有这三条性质 (i)、(ii) 与 (iii) 的开集族 \((\mathbf{U}(t))\) \((0 \leqslant t \leqslant 1)\) 。对每个 \(x \in \mathbf{X}\) ，令 \(g(x) = 1\) （若 \(x \in \mathbb{C}\mathrm{U}(\mathrm{I})\) ）；若 \(x \in \mathrm{U}(\mathrm{I})\) ，令 \(g(x)\) 为诸值 \(t\) （使 \(x \in \mathrm{U}(t)\) 者）的下确界。显然 \(0 \leqslant g(x) \leqslant 1\) （对每个 \(x \in \mathrm{X}\) ），\(g(x) = 0\) 在 \(\mathrm{A}\) 上，\(g(x) = 1\) 在 \(\mathrm{B}\) 上；而且 \(g\) 在 \(\mathrm{X}\) 上连续，因为若令 \(g(x) = a\) ，则 \(|g(y) - g(x)| \leqslant \varepsilon\) 对一切 \(y \in \mathrm{U}(a + \varepsilon) \cap \mathbb{C}\overline{\mathrm{U}}(a - \varepsilon)\) 成立——由 (I) 它是 \(x\) 的邻域——{[}约定：\(\mathrm{U}(a + \varepsilon) = \mathrm{X}\) （若 \(a + \varepsilon > 1\) ），且 \(\mathrm{U}(a - \varepsilon) = \emptyset\) （若 \(a - \varepsilon < 0\) ）{]}。

于是，只要我们能定义满足上述条件 (i)、(ii) 与 (iii) 的开集族 \((\overline{\mathbf{U}}(t))\) ，定理 1 就证明了；为此我们用公理 \((\mathbf{O}_{\mathbf{V}}^{\prime \prime})\) 。取 \(\mathbf{U}(\mathbf{i}) = \mathbb{C}\mathbf{B}\) ；由于 \(\mathbf{A} \subset \mathbf{U}(\mathbf{i})\) ，存在开集 \(\mathbf{U}(\mathbf{o})\) 使得 \(\mathbf{A} \subset \mathbf{U}(\mathbf{o})\) 且 \(\mathbf{U}(\mathbf{o}) \subset \mathbf{U}(\mathbf{i})\) （由 \((\mathbf{O}_{\mathbf{V}}^{\prime \prime})\) ）。设对每个二进数 \(k / 2^n\) （ \(k = \mathbf{o}, \mathbf{i}, \ldots, 2^n\) ）已定义了开集 \(\mathbf{U}(k / 2^n)\) ，这些集使得 \(\overline{\mathbf{U}}(k / 2^n) \subset \mathbf{U}((k + \mathbf{i}) / 2)^n\) （对 \(\mathbf{o} \leqslant k \leqslant 2^n - \mathbf{i}\) ）。对每个二进数

\[
(2 k + \mathrm{I}) / 2 ^ {n + 1} \quad (\mathrm{o} \leqslant k \leqslant 2 ^ {n} - \mathrm{I})
\]

由 \((\mathbf{O}_{\mathbf{v}}^{\prime \prime})\) 存在开集 \(\mathrm{U}((2k + 1) / 2^{n + 1})\) 使得

\[
\overline {{{{\mathbf {U}}}}} (k / 2 ^ {n}) \subset \mathrm{U} ((2 k + \mathrm{I}) / 2 ^ {n + 1})
\]

且

\[
\overline {{{\mathbf {U}}}} ((2 k + \mathrm{I}) / 2 ^ {n + 1}) \subset \mathbf {U} ((k + \mathrm{I}) / 2 ^ {n}).
\]

于是对每个二进数 \(r\) （满足 \(0 \leqslant r \leqslant 1\) ）可定义开集 \(U(r)\) ，使得 \(A \subset U(o)\) ，\(B \subset \complement U(1)\) ，且

\[
\overline {{{\mathbf {U}}}} (r) \subset \mathbf {U} (r ^ {\prime})\tag{2}
\]

对每对二进数 \(r, r'\) （满足 \(0 \leqslant r < r' \leqslant 1\) ）成立。

现在对每个实数 \(t \in [0, 1]\) 定义

\[
\mathrm{U} (t) = \bigcup_ {r \leqslant t} \mathrm{U} (r) \quad (r \text { dyadic });
\]

由 (2)，当 \(t\) 为二进数时这个定义与前面的定义一致；又，若 \(0 \leqslant t < t' \leqslant 1\) ，则存在两个二进数 \(r, r'\) 使得 \(t \leqslant r < r' \leqslant t'\) ；由 (2) 有 \(\mathrm{U}(t) \subset \mathrm{U}(r)\) ，因此

\[
\overline {{{\mathbf {U}}}} (t) \subset \overline {{{\mathbf {U}}}} (r) \subset \mathbf {U} (r ^ {\prime}) \subset \mathbf {U} (t ^ {\prime});
\]

这就证明了 (1)，从而完成了证明。

定理 1 将使我们能够证明两类重要的拓扑空间都是正规的。首先：

命题 1. 紧空间是正规的。

由第 I 章，§ 9，第 2 号的命题 2，紧空间满足公理 \((\mathrm{O}_{\mathrm{V}}^{\prime})\) 。

至于局部紧空间，这种空间的每一点都有紧邻域，后者是正规子空间；但存在不是正规的局部紧空间的例子（参见习题 9 与 26，及 § 5 习题 15）。

命题 2. 可度量化空间是正规的。

设 X 是可度量化空间，d 是与 X 的拓扑相容的度量。设 A、B 是 X 的两个不相交闭子集；由于函数 \(d(x, \mathbf{A})\) 与 \(d(x, \mathbf{B})\) 连续，使 \(d(x, \mathbf{A}) < d(x, \mathbf{B})\) {[}相应地，\(d(x, \mathbf{B}) < d(x, \mathbf{A})\) {]} 的点 x 所成的集 U（相应地，V）是开的；显然 \(A \subset U\) 且 \(B \subset V\) 且 \(U \cap V = \varnothing\) ，故公理 \((\mathrm{O}_{\mathrm{V}}^{\prime})\) 满足。

注。1) 命题 2 给出了可度量化的又一个必要条件；但这个条件即使与 § 2 中给出的所有必要条件合在一起，也不构成可度量化的一组充分条件（参见习题 6 及 § 5 习题 10）。

\begin{enumerate}
\def\labelenumi{\arabic{enumi})}
\setcounter{enumi}{1}
\tightlist
\item
  存在既不可度量化也非局部紧的正规空间的例子（见 § 5 习题 16）。
\end{enumerate}

由 \((\mathrm{O}_{\mathrm{V}}^{\prime})\) ，正规空间的每个闭子集都是正规子空间；但对正规空间的任意子集，情况未必如此。

例如，一个非正规的完全正则空间同胚于一个紧空间的子空间（§ 1，第 5 号，命题 3），而后者是正规的。

最后我们指出，两个正规空间之积未必正规（见习题 9 及 § 5 习题 16）。

